\subsubsection{Uniqueness}\label{subsec:conservation-uniqueness}

We have now discussed several selection criteria and noted some of the relationships between them. Our stated goal was to achieve some sort of uniqueness result. After all this work, are we in a position to do this? The answer, in general, is "no." Although the criteria we have suggested rule out the most obvious "physically unreasonable" weak solutions, the question of existence and uniqueness is, in general, open. At the time of this writing, this is a very active area of research. In the following, we summarize a number of results in special cases.

For the scalar conservation law with strongly convex $f$, the questions of existence and uniqueness are basically settled. For genuinely nonlinear systems, existence (but not uniqueness) is known for initial data of small total variation. For the $p$ system, assuming strong convexity, much more is known. Solutions exist for arbitrary initial data, and uniqueness has been shown within the class of piecewise smooth solutions. We refer to~\cite{Sm} for an exposition.

There are many specialized results for other systems, e.g., those where genuine nonlinearity is violated in a specific fashion and for the system of gas dynamics. Existence results are usually based on finding estimates for approximated solutions and extracting convergent subsequences. Such approximate solutions usually come from finite difference schemes or, alternatively, from adding "viscosity" terms to the equations. Some of the main contributors to the field are Lax, Glimm, DiPerna, Tartar, Godunov, Liu, Smoller and Oleinik. Despite all of these efforts, general answers in this field have remained elusive. In fact, there are recent counterexamples where the usual admissibility conditions do not guarantee uniqueness~\cite{Se}. Of course, real world problems are usually in more than one space dimension. Almost everything is open for that situation.

\paragraph*{Problems}

\begin{enumerate}[label=\thesection.\arabic*.,ref=\thesection.\arabic*,start=1,series=book-problems]

\item\label{ex:p-system-hyperbolicity} Show that if $p'<0$, then the $p$ system is hyperbolic.

\item\label{ex:gas-dynamics-hyperbolicity} Give conditions on the constitutive functions ensuring that the two systems of gas dynamics equations are hyperbolic.

\item\label{ex:conservation-characteristics-proof} Prove Lemma~\ref{lem:conservation-characteristics}.

\item\label{ex:riccati-blowup-proof} Prove Lemma~\ref{lem:riccati-finite-time-blowup}.

\item\label{ex:p-system-blowup-proof} Prove Theorem~\ref{thm:p-system-finite-time-blowup}.

\item\label{ex:sketch-mixed-rarefaction-shock} Sketch a characteristic diagram and the wavefront for the set of solutions given by \eqref{eq:burgers-continuum-solution-data}.

\item\label{ex:convexity-lax-scalar-shock} Let $f$ be convex and let $u$ be a piecewise smooth weak solution of \eqref{eq:scalar-conservation-law-quasilinear} with a finite number of jumps. Show that if $u$ is monotone decreasing as a function of $x$, then it satisfies the Lax shock condition at each discontinuity. Use an example to show that this is false if $f$ is nonconvex.

\item\label{ex:entropy-jump-condition-proof} Prove Lemma~\ref{lem:entropy-jump-condition}.

\item\label{ex:riemann-invariant-gradient-basis} Show that the set $V$ defined in \eqref{eq:rarefaction-invariant-gradient-basis} is a basis for $\mathbb{R}^n$. Hint: What is the relationship between $\bm{r}_i$ and $\bm{l}_j$ for $i\ne j$?

\item\label{ex:p-system-shock-reparametrization} Show that the states defined by the curve defined in \eqref{eq:p-system-admissible-first-shock-states} and \eqref{eq:p-system-admissible-second-shock-states} with corresponding shock speeds satisfy \eqref{eq:p-system-back-shock-speeds} and \eqref{eq:p-system-front-shock-speeds}, respectively. Hint: Use the convexity of $p$ before taking square roots.

\item\label{ex:p-system-fast-family-covering} Show that the fast family covers a neighborhood of $\bar{\bm{u}}$ univalently.

\item\label{ex:eulerian-lagrangian-gas-equivalence} Show that Eulerian and Lagrangian gas dynamics are equivalent for smooth solutions. What are the difficulties with weak solutions?

\end{enumerate}
